\section{VR実験計画}
\label{sec:experiment}

\subsection{実験環境}

Unity（URP）上に構築した工場シーン \texttt{FactoryExperiment} を用い，
Meta Quest 3（OpenXR）で被験者に没入させる．
10〜20台のAGVが荷物のピックアップとドロップを反復する環境で，
被験者はStation AからStation Bへ移動する（\Cref{tab:layout}）．

\begin{table}[htbp]
  \centering
  \caption{実験空間の主要座標}
  \label{tab:layout}
  \begin{tabular}{ll}
    \toprule
    名称 & 座標 $(x, y, z)$ \\
    \midrule
    Station A（開始） & $(22,\, 0.2,\, 9)$ \\
    Station B（ゴール） & $(9,\, 0.2,\, 58)$ \\
    AGV床面高 & $y = \SI{0.15}{m}$ \\
    \bottomrule
  \end{tabular}
\end{table}

\subsection{被験者とデザイン}

\begin{itemize}[leftmargin=2em]
  \item \textbf{デザイン}: 被験者内計画，3条件（Baseline / No-AR / Proposed）
  \item \textbf{反復}: 各条件10ケース（計30試行／被験者）
  \item \textbf{ケース固定}: 同一ケース番号ではAGV台数（10--20），速度（\SIrange{1.0}{2.0}{m/s}），
        軌道をシード固定し，条件間でAGV挙動を一致させる
  \item \textbf{被験者数}: 予備実験後に決定（パイロット $n \ge 8$ を推奨）
\end{itemize}

\subsection{手続き}

\begin{enumerate}[leftmargin=2em]
  \item 説明と同意取得，HMD装着，操作練習
  \item 実験者がPC上のメニューから条件とケース（1--10）を選択
  \item 3秒のイントロ表示後，Station Aにテレポート
  \item 被験者はStation B（半径\SI{3}{m}以内）へ移動
  \item ゴール到着で試行終了，データ保存，Station Aへ復帰
  \item 全条件・全ケース完了まで反復
\end{enumerate}

\subsection{操作}

\begin{itemize}[leftmargin=2em]
  \item \textbf{移動}: 左スティック（HMDの向き基準）
  \item \textbf{視点}: HMDの頭部追従
  \item \textbf{歩行速度}: \SI{1.4}{m/s}
\end{itemize}

\subsection{従属変数}

\begin{table}[htbp]
  \centering
  \caption{従属変数と操作定義}
  \label{tab:dv}
  \begin{tabular}{llp{5.5cm}}
    \toprule
    変数 & 単位 & 定義 \\
    \midrule
    完了時間 & s & Station A出発からStation B到達まで \\
    ニアミス回数 & 回 & AGV半径\SI{0.5}{m} + 被験者半径\SI{0.3}{m}
    = \SI{0.8}{m} 以内接近（AGVごとクールダウン\SI{1.5}{s}） \\
    移動距離 & m & 水平面上の軌跡長（\SI{0.5}{s}間隔サンプリング） \\
    \bottomrule
  \end{tabular}
\end{table}

\subsection{データ記録}

1 PlayセッションにつきExcel（.xlsx）1ファイルを出力する．
Summaryシートに条件，ケース番号，従属変数を記録し，
ケース別シートに時系列位置 $(x,y,z)$ を保存する．
